\subsection{Cinemática Diferencial e Jacobiano}
\label{sec:cinematica_diferencial_jacobiano}

A cinemática diferencial descreve como as velocidades articulares 
$\dot{\vec q} = [\dot\theta_1,\dot\theta_2,\dot\theta_3,\dot\theta_4,\dot d_5]^T$
geram a velocidade linear e angular da ferramenta em um referencial escolhido.
Esse mapeamento é feito pela matriz Jacobiana geométrica $\mathbf{J}(\vec q)$:

\begin{equation}
\label{eq:xdot_J_qdot}
{}^{0}\!\dot{\vec x}_{E} =
\begin{bmatrix}
{}^{0}\!\vec v_{E} \\
{}^{0}\!\vec \omega_{E}
\end{bmatrix}
= \mathbf{J}(\vec q)\,\dot{\vec q}.
\end{equation}

\subsubsection{Construção do Jacobiano}

As colunas de $\mathbf{J}$ são obtidas a partir dos eixos ${}^{0}\!\vec z_{i-1}$ e das posições ${}^{0}\!\vec p$ da cinemática direta. Para uma junta revoluta $i$:

\begin{equation}
J_{v,i} = {}^{0}\!\vec z_{\,i-1} \times 
\big({}^{0}\!\vec p_{E} - {}^{0}\!\vec p_{\,i-1}\big),
\qquad
J_{\omega,i} = {}^{0}\!\vec z_{\,i-1},
\end{equation}

enquanto que, para uma junta prismática:

\begin{equation}
J_{v,i} = {}^{0}\!\vec z_{\,i-1}, 
\qquad
J_{\omega,i} = \vec 0.
\end{equation}

\subsubsection{Efeito da Base}

Se a base for considerada fixa, utiliza-se \eqref{eq:xdot_J_qdot}.
No caso de base flutuante, a velocidade da ferramenta inclui a contribuição do movimento da base, de velocidade linear e angular ${}^{0}\!\dot{\vec x}_B$:

\begin{equation}
{}^{0}\!\dot{\vec x}_{E} =
\mathbf{J}_{base}(\cdot)\,{}^{0}\!\dot{\vec x}_{B}
+ \mathbf{J}(\vec q)\,\dot{\vec q}.
\end{equation}

\subsubsection{Mapeamento de forças e torques}

O mapeamento entre a \emph{wrench} cartesiana na ferramenta e os torques de junta é dado por:

\begin{equation}
\vec\tau = \mathbf{J}^T (\vec q)\,\vec w_{E}, 
\qquad
\vec w_{E} =
\begin{bmatrix}
\vec f_{E} \\ \vec m_{E}
\end{bmatrix}.
\end{equation}

\subsubsection{Singularidades}

Singularidades ocorrem quando $\mathbf{J}(\vec q)$ perde posto, exigindo velocidades articulares elevadas para pequenas variações na ferramenta. 
Uma solução prática em controle inverso é a pseudo-inversa amortecida:

\begin{equation}
\label{eq:pseudoinversa}
\dot{\vec q} =
\mathbf{J}^T\,\big(\mathbf{J}\,\mathbf{J}^T
+ \lambda^2 I\big)^{-1}\,{}^{0}\!\dot{\vec x}_{E,d},
\end{equation}

com $\lambda > 0$ atuando como fator de regularização.
